
\hypertarget{basic-structure}{%
\subsubsection{Basic Structure}\label{basic-structure}}

There are decisive advantages to using DAGs (instead of trees) as the fundamental structure of a chain.
Namely, multiple histories (both compatible and incompatible) can be merged into a single, consistent history --- a feature which eliminates stale blocks and thwarts attacks like an empty-block DoS.
UT's simplex-chains must be block-DAGs to remain functional and avoid such DoS attacks.

Often, the motivation for using a block-DAG --- instead of a block-tree --- is to increase the block frequency.
Since block-DAGs can reference multiple parent blocks, the stale-rate can theoretically approach (or reach) 0.
In UT, increasing the block frequency eventually becomes counter-productive, though, since UT is sensitive to the size and number of headers that are produced (see \autoref{sec:impact-of-header-size}).
In UT, the purpose of using block-DAGs is to thwart certain attacks, not to increase the block frequency.
The intention is for UT simplex-chains to use fairly typical block frequencies (e.g., 1 block per 15 s) --- possibly decreasing those frequencies over time to increase capacity.
Slower block frequencies also decrease the incidence of multiple parents --- this matters because each parent typically increases the header size by 32 bytes.

Some basic block-dag segments are shown in \autoref{fig:dag-simple-segments}.

\begin{figure}[h]
    \begin{subfigure}[t]{.31\linewidth}
        \vskip 0pt
        \centering
        \includegraphics[height=.95\linewidth]{dag_simple_sag}
        \caption{A simple 2-parent example of a block-dag segment.}
    \end{subfigure}%%
    \hfill
    \begin{subfigure}[t]{.31\linewidth}
        \vskip 0pt
        \centering
        \includegraphics[height=.95\linewidth]{dag_simple2_sag}
        \caption{A slightly more complex example of a block-dag segment.}
        \label{fig:dag-simple2}
    \end{subfigure}%%
    \hfill
    \begin{subfigure}[t]{.31\linewidth}
        \vskip 0pt
        \centering
        \includegraphics[height=.95\linewidth]{dag_simple3_sag}
        \caption{A 3-parent example of a block-dag segment.}
    \end{subfigure}%%
    \caption{Some examples of simple block-dag segments.}
    \label{fig:dag-simple-segments}
\end{figure}

The essence of DAG-based consensus (at least the kind we're concerned with) is to \emph{prioritize execution} of blocks and transactions based on the \emph{security contribution} (i.e., weight) of each parent block (and that parent's ancestry).
Based on a \emph{most recent common primary\footnotemark ancestor}, we can decide which parent's \emph{history} has priority execution.
Prioritized blocks are positioned \emph{earlier} in the final ordering.

\footnotetext{
    In DAG (or DAG-like) chains, primary ancestors form the \emph{main chain}, aka. the \emph{pivot chain}.
    Provided that parent blocks \emph{are prioritized by cumulative work}, the chain of \emph{primary ancestors} between a given block and the genesis block must be the \emph{heaviest path} between the two.
}

When a block has more than one parent, the prioritized parent is the \emph{primary} one.
The chain of primary parents forms the \emph{main chain}.

\defineTermTex{Main Chain}{%
    In a block-DAG, the \emph{main chain} is the continuous chain of primary parents from the best block to the genesis block.
    Blocks that are part of the main chain are \emph{on-main}, and blocks that are not are \emph{off-main}%
}

Let's limit block-DAG blocks to two parents.
If the best block has two parents, then each parent will have a \emph{subgraph of blocks} between itself and the \emph{most recent common primary ancestor} of the two parents.
The subgraph which takes priority is that of the \emph{prioritized parent's ancestry}.
If that subgraph is a chain, then the ordering and execution of blocks is trivial.
If it is not, then there must be another subgraph within that subgraph, and this algorithm is applied recursively to the prioritized parent.
After the prioritized subgraph is processed, the remaining blocks (those that are only ancestors of the secondary parent block, if it exists) are ordered and applied --- invoking recursion where necessary.
Finally, the best block is applied.
In this way, all blocks are executed after their ancestors, and there is a clear and total ordering that trivially converges.

There is a simple generalization of the above to allow more than two parent blocks, too.
That is: replace \emph{all} but the last (worst) parent with a \emph{virtual parent block} that links back to all remaining \emph{actual} parent blocks (but contributes zero block-weight itself).
Replacing the best parents (rather than the worst parents) with a virtual block means that the fork rule works when selecting the virtual block over the least-prioritized parent.
This can be repeated to allow for arbitrarily many parents.

\autoref{fig:dag-ex1-full} is an example of this algorithm for a moderately complex chain-segment (\(\BB\cdots \BB{3}\) which is 7 blocks total), and each step is enumerated and explained.

\def\dagExamplesWidth{.9375}
\begin{figure}[p]
    \begin{subfigure}[t]{.32\textwidth}
        \vskip 0pt
        \centering
        \includegraphics[width=\dagExamplesWidth\linewidth]{dag_example1_sag}
        \vspace{0.55em}
        \caption{
            How should we order this block-DAG?
            Note: children should \emph{always} be after their parents, and \emph{prioritization} means \emph{earlier execution}.
        }
        \label{fig:dag-ex1}
    \end{subfigure}%%
    \hfill
    \begin{subfigure}[t]{.32\textwidth}
        \vskip 0pt
        \centering
        \includegraphics[width=\dagExamplesWidth\linewidth]{dag_example1_expanded_order_00_sag}
        \caption{The first thing we should do is create any virtual nodes that are required ($V_{i+2,1}$).}
        \label{fig:dag-ex1-order-00}
    \end{subfigure}%%
    \hfill
    \begin{subfigure}[t]{.32\textwidth}
        \vskip 0pt
        \centering
        \includegraphics[width=\dagExamplesWidth\linewidth]{dag_example1_expanded_order_10_sag}
        \caption{Since there are two parents, we look at the \emph{prioritized subgraph} (i.e., most worked).}
        \label{fig:dag-ex1-order-10}
    \end{subfigure}

    \begin{subfigure}[t]{.32\textwidth}
        \vskip 0pt
        \centering
        \includegraphics[width=\dagExamplesWidth\linewidth]{dag_example1_expanded_order_20_sag}
        \caption{Again, there are two parents, so we look at the \emph{next} prioritized subgraph.}
        \label{fig:dag-ex1-order-20}
    \end{subfigure}%%
    \hfill
    \begin{subfigure}[t]{.32\textwidth}
        \vskip 0pt
        \centering
        \includegraphics[width=\dagExamplesWidth\linewidth]{dag_example1_expanded_order_30_sag}
        \caption{
            We've found a chain.
            These blocks have the highest priority, so are executed first.
            }
        \label{fig:dag-ex1-order-30}
    \end{subfigure}%%
    \hfill
    \begin{subfigure}[t]{.32\textwidth}
        \vskip 0pt
        \centering
        \includegraphics[width=\dagExamplesWidth\linewidth]{dag_example1_expanded_order_40_sag}
        \caption{
            We're now \emph{ordering} the blocks --- the solid arrows represent the final ordering.
            In this step we order the highest priority blocks.}
        \label{fig:dag-ex1-order-40}
    \end{subfigure}

    \begin{subfigure}[t]{.32\textwidth}
        \vskip 0pt
        \centering
        \includegraphics[width=\dagExamplesWidth\linewidth]{dag_example1_expanded_order_50_sag}
        \caption{Now that the highest priority blocks are ordered, we can order the \emph{previous} subgraph.}
        \label{fig:dag-ex1-order-50}
    \end{subfigure}%%
    \hfill
    \begin{subfigure}[t]{.32\textwidth}
        \vskip 0pt
        \centering
        \includegraphics[width=\dagExamplesWidth\linewidth]{dag_example1_expanded_order_60_sag}
        \caption{Once more, we order the next-in-line subgraph.}
        \label{fig:dag-ex1-order-60}
    \end{subfigure}%%
    \hfill
    \begin{subfigure}[t]{.32\textwidth}
        \vskip 0pt
        \centering
        \includegraphics[width=\dagExamplesWidth\linewidth]{dag_example1_expanded_order_70_sag}
        \caption{And finally we order the last remaining blocks.
            (We could remove virtual blocks too).}
        \label{fig:dag-ex1-order-70}
    \end{subfigure}
    \caption[
        Example: sorting a moderately complex block-DAG.
    ]{
        Example: sorting a moderately complex block-DAG; note that the left parent is always the best parent, so will have priority.
        Each block is annotated with its \emph{chain-weight} ($\Sigma_w$).
    }
    \label{fig:dag-ex1-full}
\end{figure}

We should expect conflicting transactions (which might otherwise be attempted doublespends) to arise during this process.
Ancestors of one parent may not be ancestors of another parent.
The exact protocol for handling conflicts is up to the implementation, but there is no reason that secondary parents should be treated as invalid by default.
We will discuss merging histories in \autoref{sec:merging-histories}.

Further reading: \citeInclusiveFull{}.
